\documentclass[aspectratio=169,11pt]{beamer}

\usetheme{Madrid}
\usecolortheme{default}
\usefonttheme{professionalfonts}
\setbeamertemplate{navigation symbols}{}
\setbeamertemplate{footline}[frame number]

\usepackage{amsmath, amssymb, booktabs}

\title{Matching, Market Design, And Labor Allocation}
\subtitle{Centralized Matching As A Labor-Market Institution}
\author{Special Topic 7: Labor Market Design, Contracting, and Mechanism Design -- Week 1}
\date{}

\begin{document}

\begin{frame}
  \titlepage
\end{frame}

\begin{frame}{Central question}
\begin{itemize}
    \item When do designed matching rules improve labor allocation relative to decentralized search and bargaining?
    \item The week is theory-heavy, but the theory is anchored in labor-market objects: jobs, training slots, timing, priorities, career starts, staffing, and welfare.
    \item The research standard is translation: institution $\rightarrow$ design failure $\rightarrow$ data $\rightarrow$ counterfactual $\rightarrow$ welfare or fairness.
\end{itemize}
\end{frame}

\begin{frame}{Why centralized matching appears in labor markets}
\begin{itemize}
    \item Professional entry markets often combine thick participation with slow offer-response cycles.
    \item Decentralized offers can create early contracting, bottlenecks, exploding deadlines, and unstable assignments.
    \item A clearinghouse pools applicants and programs at one decision moment.
    \item Centralization shifts competition from timing to preferences, priorities, capacities, and mechanism rules.
    \item The labor question is comparative: which failure does the design solve, and what new incentives does it create?
\end{itemize}
\end{frame}

\begin{frame}{Stability, thickness, congestion, and timing}
\small
\begin{tabular}{p{0.22\linewidth} p{0.33\linewidth} p{0.32\linewidth}}
\toprule
Object & Labor-market meaning & Empirical signal \\
\midrule
Stability & no applicant-program pair wants to defect & rematching pressure, off-cycle deals \\
Thickness & participants arrive together & enough comparable options at one date \\
Congestion & offers cannot clear smoothly & queues, delays, rejected offers, vacancies \\
Timing & offers move before information is ready & early acceptances, exploding offers \\
\bottomrule
\end{tabular}
\vspace{0.25cm}
\begin{itemize}
    \item These are not decorative concepts; they name what the empirical design must observe.
\end{itemize}
\end{frame}

\begin{frame}{Strategy-proofness, priorities, and fairness}
\begin{itemize}
    \item Strategy-proofness asks whether truthful rankings are a best response for one side.
    \item Priorities determine whose claim to scarce training slots is honored by the mechanism.
    \item Fairness is not automatic: a stable assignment can still redistribute opportunity or leave staffing goals unmet.
    \item Labor-market priorities may encode couples, specialties, geographic needs, diversity goals, service obligations, or public staffing constraints.
    \item Welfare requires a named object: match quality, access, strategic burden, staffing, wages, or worker welfare.
\end{itemize}
\end{frame}

\begin{frame}{Professional-entry labor markets as the anchor}
\small
\begin{tabular}{p{0.25\linewidth} p{0.34\linewidth} p{0.29\linewidth}}
\toprule
Market & Design problem & Labor object \\
\midrule
Medical residency & stable many-to-one matching, couples, timing & career starts and hospital staffing \\
Clinical psychology & bottlenecks and delayed information & congestion and early contracting \\
Specialist fellowships & unraveling and thin submarkets & timing rules and mobility \\
Law clerkships & informal norms and exploding offers & prestige, bargaining, and careers \\
Public-service placement & priorities and geographic staffing & access and underserved areas \\
\bottomrule
\end{tabular}
\end{frame}

\begin{frame}{From theory to empirical design}
\begin{enumerate}
    \item Institutional detail: who ranks, who has priority, when offers arrive, and whether wages are fixed.
    \item Design object: instability, congestion, unraveling, strategic ranking, staffing, or unfair access.
    \item Data: preferences, priorities, matches, timing, vacancies, wages, locations, or outcomes.
    \item Counterfactual: decentralized offers, new timing, altered priorities, quotas, subsidies, or capacities.
    \item Welfare or fairness: match quality, strategic burden, geographic access, staffing, wages, or worker welfare.
\end{enumerate}
\end{frame}

\begin{frame}{Identification and policy comparison}
\begin{itemize}
    \item Case evidence can show why a redesign was needed and which constraints shaped implementation.
    \item Rule changes can compare outcomes before and after centralized matching, timing reforms, or priority changes.
    \item Structural matching models use observed assignments and rank data to estimate preferences and simulate counterfactual policies.
    \item The main threats are endogenous rule adoption, incomplete preference data, equilibrium response, and welfare objects hidden from the match record.
\end{itemize}
\end{frame}

\begin{frame}{Research Lab design}
\begin{itemize}
    \item Primary anchor: Roth--Peranson redesign of the medical residency match.
    \item Challenge anchor: Agarwal-style empirical and policy analysis of the medical match.
    \item Reproduce: run a simplified centralized many-to-one match using synthetic applicant-program data.
    \item Diagnose: compare centralized matching to stylized decentralized exploding offers.
    \item Transfer: redesign the workflow for clinical psychology, law clerkships, fellowships, or public-service placement.
\end{itemize}
\end{frame}

\begin{frame}{Takeaways}
\begin{itemize}
    \item Centralized matching appears when timing, congestion, and strategic behavior make decentralized hiring perform poorly.
    \item Stability, strategy-proofness, priorities, fairness, and welfare are labor-market design objects, not just theory vocabulary.
    \item Professional entry markets make the labor stakes visible: careers, training, staffing, geography, wages, and worker welfare.
    \item The frontier contribution is often empirical: link the institutional rule to a measurable allocation problem and a credible counterfactual.
\end{itemize}
\end{frame}

\end{document}
